\documentclass[11pt,a4paper]{article}
\usepackage[margin=19mm,headheight=15pt]{geometry}
\usepackage{fontspec}
\setmainfont{Noto Serif}
\usepackage{amsmath,amssymb,mathtools}
\usepackage{xcolor}
\usepackage{fancyhdr}
\usepackage{enumitem}
\usepackage{microtype}
\setlength{\parindent}{0pt}
\setlength{\parskip}{6pt}
\definecolor{navy}{HTML}{183153}
\definecolor{lightgray}{HTML}{F3F5F7}
\pagestyle{fancy}\fancyhf{}
\lhead{\small ME 5309 -- Mid-Semester Examination}
\rhead{\small ME 5309}
\cfoot{\thepage}
\begin{document}
\begin{center}
{\Large\bfseries ME 5309 -- Mid-Semester Examination}\\[3pt]
{\large Answer Booklet}\\[8pt]
\end{center}
\textbf{Instructions:} Attempt any three questions. Each question carries 10 marks. Show working where appropriate.
\section*{Question 1 \hfill [10 marks]}
\subsection*{Q1(a) \hfill [4 marks]}

(i) f\_x = 16/0.005 = 3200 pixels.
(ii) Sensor width = 6.4 mm.
I used W = 2Z tan(f/(2Z)) and obtained about 16 mm; pixel scale = 16/1280 = 0.0125 mm/pixel.

\subsection*{Q1(b) \hfill [3 marks]}

The smallest resolvable force is 0.122 N.

\subsection*{Q1(c) \hfill [3 marks]}

Using forward Euler, F approximately 1 + AT = 0.8 and G approximately BT = 0.1.
Hence x[k+1] approximately 0.8x[k] + 0.1u[k].

\section*{Question 2 \hfill [10 marks]}
\subsection*{Q2(a) \hfill [2 marks]}

At equilibrium the velocity should be zero.
The torque and gravity torque must balance.

\subsection*{Q2(b) \hfill [2 marks]}

\[\begin{aligned}A = [[0,1], [-61.3125,-2.5]], B = [[0],[50]]\end{aligned}\]

\subsection*{Q2(c) \hfill [4 marks]}

The balance is 50(90-θss) = 2.4525 sin θss.
Solving gives a small error of about 0.05°.

\subsection*{Q2(d) \hfill [2 marks]}

Gravity compensation makes the steady-state error zero.
It does not change the closed-loop poles because it only adds a constant torque.

\section*{Question 4 \hfill [10 marks]}
\subsection*{Q4(a) \hfill [2 marks]}

Closed-loop characteristic polynomial is s²+(a+bk₂)s+(c+bk₁)=s²+4s+8, giving poles -2±2j.
At DC, θ/r = bNff/(c+bk₁), so unity gain requires Nff=(c+bk₁)/b=(6+2)/1=8.\par\smallskip Comparing coefficients with $s^2+4s+8$ gives $a+bk_2=4$ and $c+bk_1=8$. The feedforward gain is chosen from the nominal plant so the reference scaling is unity; here this gives $N_{ff}=8$. The closed-loop roots of $s^2+4s+8$ are $-2\pm2j$.

\subsection*{Q4(b) \hfill [2 marks]}

With actual b=0.75 and unchanged Nff=8, k₁=2, the steady-state relation is θss = bNff/(c+bk₁) = 0.75x8/(6+0.75x2)=6/7.5=0.8.
For r=1, percentage error = (1-0.8)x100\%=20\%.\par\smallskip With the actual gain, the steady-state output is $y_{ss}=b_{actual}N_{ff}r/ (c+b_{actual}k_1)$ for the stated setup. Substituting the given numbers gives $y_{ss}=0.8$ for a unit reference. Therefore $e_{ss}=1-0.8=0.2$, i.e. a 20\% error.

\subsection*{Q4(c) \hfill [4 marks]}

The characteristic polynomial is s³+(a+bk₂)s²+(c+bk₁)s+bK\_I.
I use the target s³+6s²+12s+16, giving k₂=4, k₁=6, K\_I=16.

\subsection*{Q4(d) \hfill [2 marks]}

\vspace{2mm}

\end{document}